\begin{exercisechunk}
\item \textbf{Random examples, memory, and teacher-chosen examples.}\suppmark
\label[exercise]{ex:l8-parity-memory}
\exerciseprofile{3}{2}{3}{\exproof\quad\excalculation\quad\exconstruction\quad\exinterpretation}{%
Identify the roles of working memory, target uncertainty, and control of
examples in a parity-learning lower bound.}
An unknown set $S\subseteq[d]$ has size $k\le d/2$. A labeled example is
$(U,V)\in\mathbb F_2^d\times\mathbb F_2$, with
\[
V=\sum_{j\in S}U_j \pmod2.
\]
A streaming learner processes the examples in order, keeping a state of
at most $M$ bits between examples. In a $q$-pass algorithm it traverses
the same sequence $q$ times in the same order.

You may use the following consequence of the constant-pass learning theorem
without proof. Fix $q\in\mathbb N$. Choose $S$ uniformly among the
$k$-subsets, and independently choose the input examples uniformly and
independently from $\mathbb F_2^d$. If $k\to\infty$ and $M=o(dk)$, then
achieving $\P(\widehat S=S)\ge2/3$ requires $2^{\Omega(k)}$ examples.
The constants may depend on $q$.

\begin{enumerate}[(a),beginpenalty=10000]
\item Suppose all state retained by a training algorithm can be stored in
$O(D^2+Dm)$ numbers of $O(\log d)$ bits each. Set
\[
k=\left\lfloor\frac d{(\log d)^4}\right\rfloor,\qquad
D=\left\lceil k(\log d)^{11/10}\right\rceil,\qquad m=2k.
\]
As $d\to\infty$, show that the memory bound is $o(dk)$ and that the
theorem requires more than any polynomial number of random examples to
recover $S$ with probability at least $2/3$ in a fixed number of passes.

\item A linear-system learner can store $n=d+2$ examples in an augmented
matrix over $\mathbb F_2$. State its storage in bits and explain why its
polynomial sample guarantee is consistent with the stated lower bound.
You may use the Lecture 7 guarantee
$\P(\widehat S\ne S)<2^{d-n}$ and an in-place linear-system solver.
Next, construct a teacher-chosen set of $d$ inputs that permits exact
recovery of every $S$ while retaining only $O(d)$ bits. The teacher
chooses the inputs before seeing any labels.

\item Compare the supplied memory theorem with the fixed-parity result
in \cref{thm:classification}. State, for each, the quantifier on $S$ and
the restrictions on the learner. Explain why the memory theorem permits
the teacher-chosen protocol in (b), and why it cannot establish failure
for a learner told $S$ in advance.
\end{enumerate}
\begin{hints}
For (a), divide $(D^2+Dm)\log d$ by $dk$.
For (b), consider the standard basis of $\mathbb F_2^d$.
For (c), distinguish information supplied to the learner from knowledge
held by the person choosing the task.
\end{hints}
\end{exercisechunk}
